%% 03-data.tex — Data provenance

\section{Data Provenance}
\label{sec:data}

All input data are publicly available from the U.S. Census Bureau
at no cost.
No proprietary data, no subscription services, no restricted files.

\subsection*{Census Redistricting Files (PL 94-171)}

The primary data source is the Public Law 94-171 redistricting file,
released by the Census Bureau after each decennial census.
These files provide population counts at the census-tract and census-block
level, including total population and voting-age population by race and
Hispanic origin.

\begin{table}[h]
\centering
\caption{PL 94-171 data files used in this portfolio.}
\label{tab:pl94}
\begin{tabular}{@{}llll@{}}
\toprule
Year & Release date & Release ID & SHA-256 (placeholder) \\
\midrule
2000 & 2001-03-01 & PL-94-171-2000 & \texttt{a1b2c3d4...} \\
2010 & 2011-02-03 & PL-94-171-2010 & \texttt{e5f6a7b8...} \\
2020 & 2021-08-12 & PL-94-171-2020 & \texttt{c9d0e1f2...} \\
\bottomrule
\end{tabular}
\end{table}

\noindent The placeholder SHA-256 values above will be replaced with
confirmed hashes after the first clean data fetch using
\texttt{redist fetch --year 2020}.
Verification of individual file hashes is performed by
\texttt{redist label-verify} (see Section~\ref{sec:verification}).

\noindent \textbf{Download URLs} (automated via \texttt{redist fetch}):
\begin{itemize}[itemsep=2pt]
  \item 2020: \url{https://www2.census.gov/programs-surveys/decennial/2020/data/01-Redistricting_File--PL_94-171/}
  \item 2010: \url{https://www2.census.gov/census_2010/01-Redistricting_File--PL_94-171/}
  \item 2000: \url{https://www2.census.gov/census_2000/datasets/redistricting_file--pl_94-171/}
\end{itemize}

\subsection*{TIGER/Line Shapefiles}

Geographic boundaries for census tracts are obtained from the
Census Bureau's TIGER/Line shapefile series~\cite{tiger2020}.
These files provide tract polygons, from which the pipeline computes
shared boundary lengths used as edge weights.

\begin{itemize}[itemsep=2pt]
  \item Download path: \url{https://www.census.gov/geographies/mapping-files/time-series/geo/tiger-line-file.html}
  \item File pattern: \texttt{tl\_\{year\}\_\{fips\}\_tract.zip} (one per state per year)
  \item Total download size: approximately 40 GB for all 50 states, all three years
\end{itemize}

\subsection*{Fetching Data with \texttt{redist fetch}}

The recommended way to obtain data is through the \texttt{redist fetch} command,
which downloads, validates, and caches all required files:

\begin{verbatim}
redist fetch --year 2020               # Download 2020 data (all 50 states)
redist fetch --year 2020 --state VT    # Vermont only (~2 min, good for testing)
redist fetch --year 2020 --check-only  # Verify cache, do not download
\end{verbatim}

\noindent The fetch command validates the SHA-256 of every downloaded file
against the manifest embedded in the \texttt{redist} binary.
If a file is corrupted or modified, the command refuses to proceed.

\subsection*{Election Data}

Partisan data used in partisan-outcome metrics (Papers C.5, B.11, and related analyses)
uses the Harvard--Fekrazad 2020 tract-level election dataset, a precinct-to-tract
crosswalk derived from official 2020 general election returns.
The dataset is downloaded and processed via:

\begin{verbatim}
python scripts/data/elections/fetch_elections.py fetch --year 2020
\end{verbatim}

\noindent This step is separate from \texttt{redist fetch} and must be run independently.
SHA-256 of the downloaded election file: \texttt{[placeholder --- run \texttt{sha256sum}
on the downloaded file after \texttt{fetch\_elections.py} completes]}.

\noindent All SHA-256 hashes in this document use standard UTF-8 encoding of the
JSON canonical form.

\subsection*{Data License}

All U.S. Census Bureau data are in the public domain.
No restrictions apply to their use, reproduction, or redistribution.
See \url{https://www.census.gov/about/policies/privacy/data-linkage/privacy-confidentiality.html}.
